\chapter*{Symbols}
Symbols The following list shows the basic associations between symbols and their meanings in this book. The list is not comprehensive. Subscripts and superscripts refine the basic meanings. For example, vJ = velocity across a joint; ag = gravitational acceleration; Ii = inertia of body i; and so on. Leading superscripts denote coordinate systems. Many symbols have more than one meaning, and a few meanings have more than one symbol. ×∗, X∗dual (i.e., force-vector version of) cross-product operator (×) and coordinate transform (X) α, ˙α vectors of velocity and acceleration variables, especially as an alternative to ˙q and ¨q γ explicit motion constraint function (q = γ(y)) δ loop-closure position error; vector with one nonzero element ζ, ˙ζ vectors of contact-separation velocity and acceleration variables ι spatial impulse vector λ(i), µ(i), κ(i), ν(i) parent of (λ), set of children of (µ), set of joints that support (κ), and set of bodies in subtree starting at (ν) body i (§4.1.4) λ vector of constraint force variables; vector of contact normal forces σ spatial bias velocity in (e.g.) rheonomic constraint due to explicit dependence on time τ vector of joint-space or generalized force variables; active force variables for a joint (τ = STfJ) φ implicit motion constraint function (φ(q) = 0) Φ spatial inverse inertia ω 3D angular velocity vector a spatial acceleration vector b spatial bias acceleration (as in a = Φf + b) c velocity-dependent spatial acceleration term; 3D vector locating centre of mass C joint-space or generalized bias force (as in τ = H ¨q + C) d, e (Pl¨ucker) basis vectors on motion (d) and force (e) spaces

E 3D rotation matrix f 3D linear force; spatial force (sometimes ˆf); general element of F6 G, g coefficients of explicit motion constraint equation ( ˙q = G ˙y, ¨q = G¨y+g where G = ∂γ/∂y (if γ defined), g = ˙G ˙y) h spatial momentum vector; first moment of mass (h = mc) H joint-space or generalized inertia ¯I 3D rotational inertia I spatial inertia (IA, Ia: articulated-body, Ic: composite-rigid-body) J, JL body (J) and loop (JL) Jacobian: v = J ˙q and vJ = JL ˙q, where v and vJ are body and loop-joint velocities, respectively K, k coefficients of implicit motion constraint equation (K ˙q = 0, K¨q = k; K = ∂φ/∂q, k = −˙K ˙q) m general spatial motion vector (any element of M6) Mn, Fn vector spaces of generalized (n=n), spatial (n=6) and planar (n=3) motion (M) and force (F) vectors N, n number of things (N) and variables (n); e.g., NB, NJ, NL, n: number of bodies, joints, kinematic loops and degrees of motion freedom n moment or couple; spatial contact normal N collection of contact normals: N = [n1 n2 · · · ] O, P points in space; origins of coordinate frames; names of coordinate systems having those origins p(i), s(i) predecessor (p) and successor (s) bodies of joint i p spatial bias force (as in f = Ia + p or f = IAa + pA) P system-wide mapping between vectors pertaining to bodies and vectors pertaining to joints (6NB × 6NJ matrix) q, ˙q, ¨q vectors of joint-space or generalized position, velocity and acceleration variables; such variables for a single joint r 3D position vector of a point (e.g. an origin or a point in a body) S, T general vector subspaces; motion freedom (S) and constraint and active force (T and Ta) subspaces S, T matrices representing S and T; (T only) collection of contact normals: T = [t1 t2 · · · ] t joint-space or generalized contact normal u general vector; generalized force corresponding to ˙y; actuated subset of τ; generalized impulse v general vector; 3D linear velocity; spatial velocity (sometimes ˆv) X coordinate transform matrix (jXi: transform from i to j coordinates; XJ: joint transform; XT, XP and XS: see §4.2) y, ˙y, ¨y independent variables in explicit constraint equations; independent subset of variables q, ˙q and ¨q after application of constraint
